% ============================================================
% §4  Dynamic Fundamental Constants
% ============================================================
\section{Dynamic fundamental constants}\label{sec:stale}

In standard physics $c$, $\hbar$, and $G$ are universal constants.
In TGP, where space is generated by matter, these ``constants'' become
\textbf{fields dependent on the local density of generated space~$\Phi$}.
The only truly constant quantity is the Planck length~$\lP$.

% ----------------------------------------------------------
\subsection{Rationale: constants describe propagation}\label{ssec:uzasadnienie}

$c$ sets the speed of light propagation,
$\hbar$~the scale of quantum effects,
$G$~the strength of gravitational interaction.
All three characterise properties of the \emph{space in which propagation occurs}.
Since space in TGP is not a fixed background but a dynamical field~$\Phi$,
it is natural that the propagation properties---and hence the values
of $c$, $\hbar$, $G$---depend on the local state of that field.

% ----------------------------------------------------------
\subsection{Speed of light $c(\Phi)$}\label{ssec:c}

\begin{axiom}[Speed of light as a field]\label{ax:c}
The speed of light is a local field dependent on~$\Phi$:
\begin{equation}\label{eq:cPhi}
    c(\Phi) = c_0 \left(\frac{\PhiZero}{\Phi}\right)^{1/2},
\end{equation}
where $c_0$ is the observed speed of light at the reference
density $\Phi = \PhiZero$.
\end{axiom}

Interpretation:
\begin{itemize}[nosep]
  \item The denser the generated space (larger~$\Phi$),
        the \textbf{slower} light propagates.
  \item $\Phi \to \infty$: $c \to 0$---propagation ceases
        (``frozen'' space; black-hole interior).
  \item $\Phi \to 0$: $c \to \infty$---unphysical, since
        $\Phi = 0$ means the absence of space.
\end{itemize}

Dense space slows light---analogous to a material medium with
a large refractive index, except that here no medium exists:
the \emph{fabric of space itself} is denser.

\begin{remark}[Interpretation of $c(\Phi)$: local physical
speed]\label{rem:c-interpretacja}
In standard GR the locally measured speed of light is
\emph{always}~$c_0$ (equivalence principle); variability of~$c$
appears only when \textbf{comparing different regions}.

In TGP, Axiom~\ref{ax:c} defines $c(\Phi)$ as the
\textbf{local physical speed of light}---the speed following
from the effective metric
$g_{\mu\nu}^{\mathrm{eff}}(\Phi)$ at a given point:
\begin{equation}\label{eq:c-lok}
    c_{\mathrm{lok}}(\Phi)
    = c_0\sqrt{\frac{\PhiZero}{\Phi}}.
\end{equation}
An \emph{operationally local} observer, using instruments
made of matter immersed in the same~$\Phi$, measures~$c_0$
(self-consistency), because all measurement standards
($c$,~$\hbar$,~$G$) scale coherently with~$\Phi$
in a manner that preserves~$\lP$.
The variability of $c(\Phi)$ is therefore \emph{relational}:
it is revealed only when comparing regions of different~$\Phi$.

This distinguishes TGP from standard GR, where the variable
coordinate speed is an artefact of the coordinate choice.
In TGP $c(\Phi)$ is a \textbf{physical} consequence of
differing densities of generated space.  The difference is
testable: in TGP the ratio $c(\Phi_1)/c(\Phi_2)$ is
\textbf{observable} and depends on $\Phi_1/\Phi_2$;
in GR the analogous ratio depends on \emph{curvature},
not on any scalar field.
\end{remark}

% ----------------------------------------------------------
\subsection{Planck constant $\hbar(\Phi)$}\label{ssec:hbar}

\begin{axiom}[Emergent Planck constant]\label{ax:hbar}
The effective Planck constant is a field:
\begin{equation}\label{eq:hbarPhi}
    \hbar(\Phi) = \hbar_0 \left(\frac{\PhiZero}{\Phi}\right)^{1/2},
\end{equation}
where $\hbar_0$ is the observed Planck constant.
\end{axiom}

Interpretation:
\begin{itemize}[nosep]
  \item Dense space ($\Phi \gg \PhiZero$):
        $\hbar \to 0$---quantum effects are suppressed.
        Matter behaves classically even at the level of
        individual particles.
  \item Dilute space ($\Phi \sim \PhiZero$):
        $\hbar \sim \hbar_0$---standard quantum mechanics.
\end{itemize}

This yields \textbf{two paths to the classical limit}:
\begin{enumerate}[nosep]
  \item \emph{Statistical path}: $N \to \infty$
        (many particles; relative fluctuations vanish).
  \item \emph{Spatial path}: $\Phi \to \infty$
        (dense space suppresses quantumness).
\end{enumerate}

\begin{remark}[Black holes and quantumness]
Near a black hole $\Phi$ is enormous, so $\hbar \to 0$:
matter becomes \textbf{classical} in strong gravitational fields.
This is the opposite of the standard expectation (in GR quantum
effects should \emph{grow} near the singularity).  In TGP the
classical singularity is \emph{naturally removed}: matter cannot
``condense to a point'' because the surrounding space is too dense.
\end{remark}

% ----------------------------------------------------------
\subsection{Gravitational constant $G(\Phi)$}\label{ssec:G}

\begin{axiom}[Dynamic gravitational constant]\label{ax:G}
The gravitational constant is a field determined by the
condition of Planck-length constancy:
\begin{equation}\label{eq:GPhi}
    G(\Phi) = G_0 \cdot \frac{\PhiZero}{\Phi},
\end{equation}
where $G_0$ is the observed gravitational constant.
\end{axiom}

Interpretation:
\begin{itemize}[nosep]
  \item Dense space: $G$ decreases---gravity is
        \textbf{weaker}; more space ``dilutes'' the interaction.
  \item Dilute space: $G$ increases---gravity is stronger
        (but there is also less space).
\end{itemize}

% ----------------------------------------------------------
\subsection{Constancy of the Planck length}\label{ssec:lP}

The Planck length is defined as
\begin{equation}
    \lP^2 = \frac{\hbar\,G}{c^3}.
\end{equation}
Substituting the $\Phi$-dependent
expressions~\eqref{eq:cPhi},~\eqref{eq:hbarPhi},~\eqref{eq:GPhi}:
\begin{equation}
    \lP^2
    = \frac{\hbar_0\!\left(\frac{\PhiZero}{\Phi}\right)^{1/2}
            \cdot G_0\!\left(\frac{\PhiZero}{\Phi}\right)}
           {\left[c_0\!\left(\frac{\PhiZero}{\Phi}\right)^{1/2}\right]^3}
    = \frac{\hbar_0\,G_0}{c_0^3}
      \cdot \left(\frac{\PhiZero}{\Phi}\right)^{1/2+1-3/2}
    = \frac{\hbar_0\,G_0}{c_0^3}.
\end{equation}

\begin{theorem}[Constancy of the Planck scale]\label{thm:lP}
The Planck length is \textbf{invariant}:
\begin{equation}\label{eq:lP-const}
    \boxed{\lP = \sqrt{\frac{\hbar_0\,G_0}{c_0^3}} = \mathrm{const},}
\end{equation}
independently of the local value of~$\Phi$.
It is the only truly fundamental scale of TGP.
\end{theorem}

This is an \textbf{internal consistency condition}: despite
$c$, $\hbar$, $G$ being dynamical, the scale of quantum gravity
remains fixed.  The $\Phi$-dependences of
$c(\Phi)$, $\hbar(\Phi)$, $G(\Phi)$ are not independent
assumptions but \textbf{a single coherent condition}
arising from the nature of the field~$\Phi$.

% ----------------------------------------------------------
\subsection{Uniqueness of the exponents}\label{ssec:wykladniki}

The axioms above postulate specific power laws:
$c \propto \Phi^{-1/2}$,
$\hbar \propto \Phi^{-1/2}$,
$G \propto \Phi^{-1}$.
We now show that these exponents are the \textbf{unique}
ones consistent with three internal requirements of TGP.

\begin{theorem}[Uniqueness of exponents]\label{thm:exponents}
Let $c(\Phi) = c_0(\PhiZero/\Phi)^a$,\;
$\hbar(\Phi) = \hbar_0(\PhiZero/\Phi)^b$,\;
$G(\Phi) = G_0(\PhiZero/\Phi)^g$.\; Assume:
\begin{enumerate}[nosep]
  \item[(W1)] \textbf{Planck-length constancy:}\;
    $\lP^2 = \hbar G/c^3 = \mathrm{const}$.
  \item[(W2)] \textbf{Speed of light from the metric:}\;
    The antipodal bridge metric
    (Proposition~\ref{prop:conformal-unique})
    yields $a = 1/2$.
  \item[(W3)] \textbf{Propagator consistency
    (universality of local physics):}\;
    The effective mass $m_{\mathrm{eff}} = mc/\hbar$
    must be $\Phi$-independent, giving $a = b$.
\end{enumerate}
Then $a = b = 1/2$ (from W2 and W3), and W1 gives
\begin{equation}
  b + g - 3a = 0
  \quad\Longrightarrow\quad
  g = 3a - b = 3\cdot\tfrac{1}{2} - \tfrac{1}{2} = 1.
\end{equation}
The \textbf{unique} solution is $(a,b,g) = (1/2,\;1/2,\;1)$---precisely
Axioms~\ref{ax:c}--\ref{ax:G}.
\end{theorem}

More generally, consider any Planck combination
$c^{\,p}\,\hbar^{\,q}\,G^{\,r}$ and ask when it is
$\Phi$-independent.  The $\Phi$-exponent is
$-\tfrac{1}{2}\,p - \tfrac{1}{2}\,q - r$.
Setting this to zero gives $p + q + 2r = 0$.
The Planck length corresponds to $(p,q,r) = (-3,1,1)$,
satisfying $-3 + 1 + 2 = 0$.  This is the
\emph{unique combination} (up to overall powers)
that is $\Phi$-invariant.

% ----------------------------------------------------------
\subsection{Consequences}\label{ssec:stale-kons}

\begin{enumerate}[nosep]
  \item \textbf{No cosmological-constant problem} in the
        standard sense: $G$ varies with~$\Phi$, so the
        ``constant'' expressed through~$G$ is not constant.
  \item \textbf{Gravitational lensing}: in regions of
        higher~$\Phi$ light travels more slowly---this
        produces lensing without invoking curvature as a
        separate concept.
  \item \textbf{Gravitational redshift}: a photon travelling
        from higher to lower~$\Phi$ accelerates ($c$~increases),
        causing a redshift---recovering the GR result.
  \item \textbf{Classical limit in strong fields}:
        near massive objects $\hbar \to 0$, removing the
        need for quantum gravity in the strong-field regime.
\end{enumerate}

% ----------------------------------------------------------
\subsection{Observability of dynamic constants}%
\label{ssec:obserwabilnosc-stalych}

Variations of $c(\Phi)$, $\hbar(\Phi)$, $G(\Phi)$ are
\textbf{relational}: a local observer always measures the
reference values $c_0$, $\hbar_0$, $G_0$, because her
measurement standards (rulers, clocks, spectrometers) scale
with~$\Phi$ in the same way.  Effects are detectable only
by \emph{comparing} two regions of different~$\Phi$.

\begin{remark}[Principle of relational observability]%
\label{rem:relacyjna}
In TGP \emph{only ratios of quantities from different spatial
regions are observable}.  There is no objective ``absolute''
value of $c$, $\hbar$, or~$G$---they are numerically equal to
$c_0$, $\hbar_0$, $G_0$ \emph{by definition} for a local observer.
TGP effects manifest as anomalies in standard scaling relations
(Hubble diagram, galaxy sizes, neutron-star lifetimes).
\end{remark}
